\section{Soloviov's Philosophical Principles of Integral Knowledge}
\subsection{Law of Development}
One of the earliest works of \textbf{Vladimir Solovyov}, published when he was 24, was \textit{The Philosophical Principles of Integral Knowledge}. There, he tackles the meaning of historical development, integral knowledge, the Absolute, and the Idea. Solovyov's doctoral thesis was a survey of \textbf{Positivism} and similar movements in which he was highly critical of \textbf{Auguste Comte}. Nevertheless, the chapter on historical development seems to owe something to Comte, perhaps in an attempt to bring Positivism into a higher understanding. In the latter part of his life, Solovyov ameliorated his opinion of Comte.

He begins by asking the question of the purpose of human existence, which is really a question about the final cause of existence, something ruled out of bounds by the modernist project. Everyone, I suppose, has asked himself that question. However, Solovyov wants to go beyond the personal to the general and asks about the goal of human existence in general. Since man is always in community, Solovyov claims that ``the dignity of the personal and immediate goals of a human life can define itself only by its relationship to that common and final goal for which it serves as the means."

The question of a goal naturally leads to the question of development, the process of achieving that goal. If the ``changing phenomena [of world events] linked with each other merely in an external manner," then clearly a common goal is a meaningless concept. Of course, that is what the Baconian project entails. Hence, progressives are always in a state of perpetual development without any articulated goal.

First and foremost, ``development presupposes a specific subject"; it is the subject that develops. First, he distinguishes between an aggregate (or heap) and an organism (or system). An aggregate cannot be said to develop. Only an organism, ``a single being that contains in itself a multitude of elements that are internally connected to each other" can develop. However, the organism exists not in isolation and the ``material of development and the impelling source of its realization come from the outside". This ``impelling source can act only in conjunction with the peculiar nature of the organism. In traditional terms, the external elements constitute the material and efficient causes of development whereas the formal and final causes are contained in the subject of development.

Change by itself, without a known point of origin that continues indefinitely, not having any specific purpose, cannot be considered to be development. Here, Solovyov defines development.

\begin{quotex}
Development is that series of immanent changes in an organic being that proceeds from a known origin and directs itself toward a known, definite goal. 

\end{quotex}
Endless development, then, is a logical absurdity. In nuce, the Law of Development contains three stages:

\begin{enumerate}
\item A known primary condition, from which it is initiated 
\item Another known condition, which is its goal 
\item A series of intermediate conditions, as a transition or means 
\end{enumerate}
We note some consequences of this. First, there is a unity in the subject, which is its Form or defining characteristic. Then, the defining characteristic is transferred to the separate components.

\paragraph{Historical Development}
Solovyov next seeks to apply the Law of Development to humankind as a real, collective organism, a task he claims had never been satisfactorily carried out. Let's be clear that this is not an obvious project. For example, libertarian philosophers, such as \textbf{Friedrich Hayek}, are consistent nominalists and hence deny the real existence of humankind or any collective. He regards the individual as a real existent, but the group as such is a convenient fiction and no more than a mere name for an aggregate of individuals.

Solovyov anticipates this objection and notes that every being and every organism possesses a necessarily collective character. He points out, for example, that the cells in your body are themselves individual entities and they exist oblivious to being a part of the human organism. Hence, the individual man can be a real, individual being while still participating in a higher system. At this point, readers may want to reconsider \textbf{Boris Mouravieff}'s description of a series of higher order worlds or systems, considered as an aspect of orthodox theology.

An organism has a static part, its components, and a dynamic part, its subsystems. The component parts of the organism of humankind are not individuals, however, but rather tribes and peoples. Solovyov will proceed to develop the common formative systems of humankind.

Like Comte, Solovyov identifies three basic forms of being: \emph{feeling}, \emph{thinking}, and \emph{active will}. Personal and subjective expressions of these forms of being cannot serve as formative principles. Only objective expressions can be part of the process of development.

\begin{itemize}
\item \textbf{Objective Beauty}. That feeling that strives to strengthen its existing condition by an objective expression of it. 
\item \textbf{Objective Truth}. That kind of thinking that strives for a specific objective content. 
\item \textbf{Objective Good}. That will that has in mind definite general goals. 
\end{itemize}

\subsection{Development of the Will}

\begin{quotex}
The Will manifests as Objective Good in three levels: prosperity, justice, and eternal bliss.
\end{quotex}

The nature of a person manifests in three forms of being: \emph{feeling}, \emph{thinking}, and \emph{willing}. \textbf{Vladimir Solovyov} is interested only in their objectifications, unlike the modern mentality which considers these forms of being as merely subjective. The first form Solovyov looks at is willing. The modern mind considers the will only insofar as it reflects personal desires or, in the abstract, as simply the \emph{will to power} for its own sake. But to regard the will as a positive principle for life, requires that it seeks its objectification, or manifestation, in definite goals, a fortiori, as an objective good.

The objectification of Will takes place in three levels, corresponding to the three castes\footnote{\url{https://gornahoor.net/?p=1647}}. In order to achieve any goal, it is first necessary to provide for the existence of persons, which depends on their relationship the external world and the activities required to extract the means for existence. This is the beginning of economic society, whose task is the organization of labour. This task belongs to the Vaishyas who are oriented to the external world. This is the material basis of society. Its primary elementary form is the family, which is based on the division of labour. We saw that the Ancient City was based on extended family relationships, up to the clan.

The next level is concerned with the relationship of people to each other and their interactions as members of a single collective whole. This is political society, or the state, whose task is the organization of workers, and this task is the responsibility of the Kshatriyas. Since there are many states, there is also the problem of international relations between states. These relations are not always pacific, so the Kshatriya caste must also include warriors. However, this does not preclude the possibility of the realization of a universal state, as in the case of the Roman Empire.

The third level is determined by the religious character of a person, or his relationship to the transcendent. Clearly, the ultimate responsibility belongs to the Brahman caste. Solovyov writes:

\begin{quotex}
A person desires not only \emph{material} existence, which is provided by economic society, and not only \emph{lawful} existence, which is given to him by political society, but he desires as well an \emph{absolute} existence, one that is complete and eternal. Only the last of these is a genuine, supreme state of well-being, the summum bonum, with respect to which material goods attained through labour, and economic and formal goods attained through political activity, serve only as the means. Since the attainment of absolute existence, or eternal and blissful life, is the highest goal for everyone alike, this goal necessarily becomes the principle of social union, which may be called spiritual or sacred society. 

\end{quotex}
These considerations may be summarized in chart \ref{tab:RoleDevelopment}:

\begin{table}[h]\centering\small
\begin{tabular}{cccc}\toprule
\textbf{Level} &
\textbf{Society} &
\textbf{Correspondence} &
\textbf{Role in development}\\\midrule
I &
Economic society &
Matter &
External basis\\
II &
Political society &
Form &
Means\\
III &
Spiritual society &
Spirit &
Goal\\\bottomrule
\end{tabular}
\label{tab:RoleDevelopment}
\caption{Role of development}
\end{table}
\begin{itemize}
\item A person must live 

Material prosperity, the fruits of his labour, do not provide a person with well-being for his essential, or True, Will. Economic activity alone is not moral.
\item A person must live with others 

To be moral, the Will must take the form of justice, which is determined by the political society, or the State. Even still, just interpersonal relationships do not produce bliss.
\item A person must relate to the transcendent 

Bliss comes from principles that are outside and beyond the natural and the human worlds.
\end{itemize}
Solovyov concludes his discussion:

\begin{quotex}
Only that society that is based directly on relations with these transcendental principles can have as its direct objective the well-being of a person in his totality and absoluteness. \emph{This is what spiritual society or the church must be.} 

\end{quotex}
Here we have a program, as yet still partial since it only concerns the Will, for any attempt at the reformation of the modern world.

\begin{itemize}
\item It must provide for the material well-being of the people, whose fundamental unit is the family, not the individual. 
\item It must provide a stable and just political system that also protects against external threats. 
\item It must be unified by a common spiritual tradition, since the relationship to the Absolute provides the ultimate meaning to man. 
\end{itemize}

\subsection{The Development of Mind}

\begin{quotex}
A person may deal with the data of conscious activity on three levels: Science, Philosophy, and Theology.
\end{quotex}

\paragraph{Science}
\textbf{Science} seeks to organize and determine the underlying laws behind the mass of factual information gleaned from the observation and experience of the external world and human life. The basic element of science is the \textbf{fact}.

Our empirical knowledge of the the physical world can never be complete. On the one hand, our senses are subject to illusions and on the other, there can never be complete certainty regarding scientific theories. They cannot be fully verified, yet are always open to being falsified\footnote{\url{https://en.wikipedia.org/wiki/Falsifiability}}. Our theories of the contingent events of the physical world are necessarily matters of opinion and hence no more than likely stories. A scientific theory is the most likely story.

However, of the three levels of knowledge, science is the easiest to master, hence, it has produced the most impressive results over the past several centuries. Science is universal in the sense that it is accessible to anyone with the proper preparation and ability to understand its mathematical models. There is no requirement for a personal transformation. That is why past assertions of an ``Aryan science" or a ``Jewish science" were misguided. Science is no respecter of persons.

This is also why science is accumulative. Each scientist can build on the work of previous scientists; there is no need to repeat all the earlier experiments once certain facts have been established. Science owes its success because it makes simplifying assumptions about the world. For example, Francis Bacon\footnote{\url{https://gornahoor.net/?p=3468}} made the methodological assumption of science to bracket out the question of formal and final causes. Hence, the scientist can focus on physical things and the relationships among them. Their function or purpose are of no concern.

Scientists often make other methodological assumptions such as the anti-philosophies of materialism, atheism, or determinism. As such, there is no objection. However, scientists, in their one-sightedness, sometimes err in claiming that these assumptions are objectively true descriptions of the world.

\paragraph{Philosophy}


\textbf{Philosophy} is the formal perfection of knowledge and its logical coherence. Its basic unit is the idea. Philosophy can only point out the necessary conditions for the genuine attainment of knowledge, but it does not provide its content. The content must come from the level below or above.

Science seeks to determine the relationships among things. The goal of philosophy is to know the Sufficient Reason\footnote{\url{http://en.wikipedia.org/wiki/Principle_of_sufficient_reason}}, that is, the reason or cause for something to exist. Philosophy develops much more slowly than Science, because each generation of philosophers must recapitulate the thoughts of previous philosophers, a tedious and demanding process.

The potential errors of science are the possibility of delusion in the observations and the possibility of a theory being falsified. The potential error of philosophy comes from starting with false first principles. In that case, all deductions from such a principle will be logically false. Another potential error comes from the temptation of the sophists, the demon of dialectics; the sophist will use the tools of logic and the vocabulary of philosophy to make a case for a predetermined position, thus having only a subjective significance. A true philosopher must be disinterested in his search for first principles and sufficient reasons.

\paragraph{Theology}
\textbf{Theology} is concerned with the universal or absolute content of knowledge. Science is concerned with particular beings, philosophy with the form of knowledge, and theology with absolute being.

Unlike Science, the truths of Theology are independent of external reality. Unlike Philosophy, the truths of Theology are above Reason. Its subject matter is the absolute first principle of all of existence. Solovyov concludes:

\begin{quotex}
Only this principle imparts genuine meaning and significance both to

\begin{itemize}
\item The concepts of Philosophy
\item The facts of Science 
\end{itemize}

Without this:

\begin{itemize}
\item The concepts of Philosophy becomes meaningless form 
\item The facts of Science become undifferentiated matter
\end{itemize}
\end{quotex}

Thus we see that genuine knowledge is \emph{knowledge of the Whole}, or the One in Neoplatonic terms, or the Absolute, the Infinite, or the Tao. This form of knowledge is Gnosis and is superior to Reason.

\subsection{Development of Feeling}

By ``feeling", \textbf{Vladimir Solovyov} is not referring to subjective and personal impressions, but rather to the higher feeling centre, insofar as it is objective and leads to creation. Again, these feelings are manifested in the material, formal, and absolute levels.

\paragraph{Technical Arts}
In material creation, the idea of beauty serves merely as an ornament in the context of utilitarian goals. Solovyov actually calls this ``technical applied art", and its highest representative is architecture. This hardly seems true in the modern world, where ugly architecture\footnote{\url{http://www.forbes.com/2002/05/03/0503home.html}} is the norm. Oddly enough, more effort is now put into making electronic gadgets aesthetic as well as functional. At this level, the expression of beauty is limited by the necessity to be functional and useful. Artisans and craftsmen are the representatives of this level.

\paragraph{Fine Art}

Fine art is a kind of creation whose defining significance comes from the aesthetic form and is expressed in purely ideal images. Fine arts come in four forms, which are organized in an ascent from matter to spirit. The fine arts are distinguished from the technical arts by not having a utilitarian necessity.

\begin{description}
\item[Sculpture ] Sculpture is the most material art, the one closest. 
\item[Painting ] Painting is more ideal than sculpture since there is no imitation of the corporeal in it. Unlike sculpture, which is three dimensional, the images in a painting are depicted on a flat surface. 
\item[Music ] Music possesses more of a spiritual quality, as it is no longer embodied in or on the material itself. Rather, it is embodied in the movement and life of its substance—sound. 
\item[Poetry ] Poetry finds its expression in the spiritual element of the human word. This is why so much of the greatest spiritual literature is in the poetic form. Poetry, along with command, was the primary way of thinking in the primordial state\footnote{\url{https://gornahoor.net/index.php?s=jaynes\&submit=Search}}. 
\end{description}

\paragraph{Mysticism}
Fine art still does not reach the ideal of Beauty in its totality, since the contents of a work of art are accidental. The content of \textbf{Absolute Beauty} must be definite, necessary, and eternal. This type of beauty is not manifested directly in the world, although objects and phenomena are reflections of it. Genuine and complete Beauty, therefore, is located in the superhuman and supernatural world.

Solovyov calls the creative relationship of feeling to the transcendental world \textbf{Mysticism}. Specifically, he means the direct spontaneous relationship of the Spirit to the transcendental world, and not the negative connotations this word carries. Solovyov lists three criteria for juxtaposing mysticism with technical arts:

\begin{enumerate}
\item Both have feeling as their basis, rather than perception and active will 
\item Both rely on imagination or fantasy, rather than reflection and external activity 
\item Both presuppose a an ecstatic inspiration, rather than a calm consciousness 
\end{enumerate}
Mysticism, for Solovyov, who himself had direct mystical experiences, is far from the popular association with the vague, unclear, and incomprehensible. He explains:

\begin{quotex}
The sphere of Mysticism possesses an absolute clarity by itself, and it alone can clarify everything else, precisely as a result of this its light is unbearable for weak and unprotected eyes, and it plunges them into absolute darkness. 

\end{quotex}
\paragraph{The Degeneration of Beauty}
Concerning the degeneration of art, Solovyov writes:

\begin{quotex}
Contemporary Western art cannot capture the ideal content of life for the simple reason that it itself has lost it. In the sphere of creative work we see in the West the same alteration in the three supremacies. In the Middle Ages unfettered art does not exist; everything is subordinated to mysticism. 

\end{quotex}
Since the Renaissance, as a result of the general weakening of religious principles, the fine art of form emerged, and art was restricted to the accidental and arbitrary. In the modern world, the technical arts prevail, which are utilitarian and practical. Art has become a trade and a business. What does it cost is the primary interest of the art collector.

Art for art's sake, for the sake of beauty, no longer exists. In the attempt to transcend form, art has instead become formless. Paintings are random, music discordant, poetry meaningless. Sculpture, as the art of three dimensions is replaced with arbitrary collections of debris and detritus. In a forced effort to appear engaged and relevant, the contemporary artist can only shock the common man rather than raise his consciousness to the Ideal of Beauty. Those who are beyond being shocked are most in need of a raised consciousness, which is an invisible target for them.

\subsection{Three Stages of Development}

\begin{quotex}
In old Europe human life received its ideal content from the Catholic faith, on the one hand, and from knightly feudalism on the other \flright{\textsc{Vladimir Solovyov}}
\end{quotex}

\textbf{Vladimir Solovyov} claims that there is a three stage process of development.

\begin{enumerate}
\item In the first stage, the levels exist in an undifferentiated form. The levels exist virtually, not actually. 
\item In the second stage, the lower levels manifest independently of the higher one and seek to exist independently. Furthermore, to assert its independence, each area rejects the corresponding areas on its own level. 
\item In the third stage, the highest level separates itself, gaining its freedom and thus bringing about the possibility of a new unity. Since none of the lower levels can actually achieve exclusive domination, they are forced to seek a higher center outside themselves. This must involve the free, conscious subordination of the lower levels to the higher as their true centre. 
\end{enumerate}
These stages will be discussed in terms of the Cosmic Cycle as well as specific historical manifestations. For reference, chart \ref{tab:SpheresofExistence} is included.

\begin{table}[h]\centering\small
\begin{tabularx}{\textwidth}{rXXX}\cmidrule(lr){2-4}
 \textbf{Spheres of} &
\textbf{Creation} &
\textbf{Knowledge} &
\textbf{Praxis}\\\midrule
\textit{Subjective Basis} &
Feeling &
Thinking &
Willing\\
\textit{Objective Principle} &
Beauty &
Truth &
Good\\
\textit{Absolute Level} &
Mysticism &
Theology &
Spiritual Society\\
\textit{Formal Level} &
Fine art &
Philosophy &
Political Society\\
\textit{Material Level} &
Technical Art &
Science &
Economic Society\\
\textit{Integration} &
Theurgy &
Theosophy &
Theocracy\\\bottomrule
\end{tabularx}
\label{tab:SpheresofExistence}
\caption{Spheres of existence}
\end{table}

\paragraph{Cosmic Cycle}

From the Traditional perspective, this development is understood this way.

\begin{enumerate}
\item In the Primordial Golden Age, there is a true unity and everything is ordered to the level of the Absolute. There is no caste distinction since everyone has this understanding. 
\item The next stage encompasses the degeneration of castes in a process of decline. The centre moves from the Absolute, to the formal, to the material following the successive revolts of the lower castes. Eventually, the material world is the main focus. However, this degeneration cannot continue unabated. 
\item In the ascending phase of the Kali Yuga, there is the possibility for some men to come to a higher awareness and adopt the higher viewpoint, not as a matter of birth, but rather as a conscious and free choice. This will lead up to a new Golden Age. 
\end{enumerate}
\paragraph{Historical Manifestation}
Solovyov does not explicitly mention the law of cosmic cycles, though he had to have known of them. However, he helpfully gives several examples of the Law of Development.

\subparagraph{Ancient City}
In the Ancient City, the spheres and levels were not differentiated. The family and clan were the first forms of economic life, but they also possessed a political and religious significance\footnote{Since we have already written extensively on this in \textit{The Mind of the Ancient City}, \textit{The Religion of the Ancient City}, and \textit{The Degeneration of the Ancient City}, we will pass over this cursorily. See its respective Sections in the \textit{Ancient City} volume.}. Solovyov notes that the levels were not distinguished and hence they formed interconnected wholes. He identifies them as:

\begin{itemize}
\item \textbf{Theosophy}: comprising theology, philosophy, and science. 
\item \textbf{Theocracy}: comprising church, government, and economic society. 
\item \textbf{Theurgy}: comprising mysticism, fine art, and technical art in a single mystical creative process. 
\end{itemize}
Together, these three form a single religious whole.

\subparagraph{Christendom}
It is difficult for many to view the period of Christendom with the necessary detachment to understand it. In late paganism, the rites and rituals of the Ancient City had become empty forms, poorly understood, if at all. It was in this milieu that a change was required. Keep in mind Guenon's claim that the initial impulse of Christianity was esoteric, but that at a certain point it was necessary for it to become exoteric. Solovyov points out the reason for this:

\begin{quotex}
It must be acknowledged that in the universal consciousness of humankind the original interconnectedness was shaken decisively and to its very foundations only with the advent of Christianity, when for the first time the sacred was separated from the profane. In this respect, as a force that dealt the final blow to this external, involuntary unity, Christianity represented the beginning of genuine freedom.

\end{quotex}
Solovyov then traces the detachment and isolation of the spheres and levels that followed in Christendom. The process begins with the separation of the two lower profane, natural levels from the sacred, divine highest level. Specifically, the government and economic activities separate from the Church. The full explication of this process will have to await a separate posting, but we will report Solovyov's conclusion.

\begin{quotex}
Christians stripped [the government] of any positive spiritual content. For them the emperor —the last god of the pagan world— could be only the chief administrator of the police. This relationship negated the very principle of ancient society, which entailed specifically the deification of the republic and the emperor as its representative, the merging of spiritual and secular principles. Hence the emperors, in persecuting the Christians, were acting not as bearers of government powers in the narrow sense of the term, but as bearers of the entire ancient consciousness.

\end{quotex}
There are many today who resent this either because they still retain elements of the ancient consciousness or, as is more likely, they resent the imposition of freedom required by the new order of things. The separation of the sacred from the profane introduces an element of choice, which interferes in the felt immediacy of a natural consciousness.

Next: The decline of Christendom and the lineaments of the next stage.

\textbf{Reformation and Revolution}

\textbf{The Third Stage}


\hfill



\flrightit{Posted on 2011-12-01 by Cologero }

\begin{center}* * *\end{center}

\begin{footnotesize}\begin{sffamily}



\texttt{logres on 2011-12-06 at 23:20 said: }

This is perfectly put together – how is it done? But it doesn't matter. A thought or a question – in a graded hierarchy (replete with caste and with shades of gradation through the human pyramid), wouldn't it actually be EASIER to relate to others, both immediately above you, and slightly beneath you? Since the difference (although perhaps great in terms of caste) would be not so great in terms of human gradation or circumstance? Barzun thought that democracy had caused a huge suspicion of intellect \& endless jockeying for position, since no one ``knew where they stood".


\hfill

\texttt{Avery on 2011-12-07 at 00:03 said: }

Hello, I'm sorry for posting my comment on your last article without reading the article itself. (It was a very weird thing to do, and I think I was in a mental pit devoid of all self-awareness.) I have now read both articles and I think the subject is fascinating. 

I agree with logres that this hierarchy is better than the status quo, but I think democracy is only part of the issue. Human beings respond to the metaphors our teachers fit us into. Economics analogizes human behavior to the physics of a machine. If our society is analogized to a machine by its smartest leaders, then we will act like we are trapped in a machine, and indeed many intellectuals seem to be ruled by this delusion. But if we are meant to be playing a role in a great Will, then we will engage ourselves to express that Will. In your model, flaws in our society come not from the rulers playing with the controls of the machine, but with everyone's personal lack of understanding of what the Will of their nation is and how they are meant to execute it.


\hfill

\texttt{Boreas on 2011-12-15 at 07:11 said: }

Good and relevant text.

``Art for art's sake, for the sake of beauty, no longer exists."

If this (for the sake of beauty) is behind ``art for art's sake", then I'm absolutely for it, but as I have understood, the (post-)moderners have twisted this concept to its opposite, when it has become to mean that art has become ``a luxury item for idle parasites" together with being transformed into a meaning in itself in the form of ugly, pointless and abstract public scenery sculptures made of…junk and rubbish.


\hfill

\texttt{Matt on 2011-12-15 at 17:33 said: }

The experience music project building looks like a collapsed lung combined with a mangled artery.


\hfill

\texttt{Charlotte Cowell on 2011-12-15 at 18:14 said: }

Does beauty alone need art or, indeed, does beauty exist outside its connection with the other attributes of God? you could call God beautiful, but could you call beauty God? You could argue that a portrait of a very beautiful woman, or a sunset, for example, were works of art that have the sole purpose of capturing beauty, but they are bound to fall short of the perfect originals and are, therefore, essentially acts of praise/worship. Beauty doesn't exist in a vacuum, for wherever you find beauty you will also find truth, power, wisdom, etc…


\hfill

\texttt{Cologero on 2011-12-15 at 20:09 said: }

I really don't understand this comment, Charlotte, whether you're disputing or expanding on a point. There is no intent to ``argue" anything, rather simply to describe. Beauty as one of the three highest ideas manifests on three planes corresponding to the tripartite constitution of man: spirit, soul, body. So it is not the case that ``beauty needs art", but rather that it manifests as art on the formal level. Solovyov makes the point I believe you are trying to make, that art at the formal level, because it always has a specific and somewhat arbitrary subject matter, is not the fullest manifestation of beauty. That is reserved to mysticism.

We have been illustrating the relationship of the Absolute with the Infinite, which manifests in all the permutations of the world process. So, of course, Beauty does not exist in a vacuum, but it is helpful to analyze how the ideas do manifest in isolation. For all practical purposes that is how we typically experience the ideas … in isolation. The scientist denies the mystical, the mystic knows nothing about economics, and so on. Wisdom is knowledge of the whole and we are here trying to gain such knowledge.


\hfill

\texttt{Charlotte Cowell on 2011-12-16 at 04:09 said: }

I wasn't doing either, just commenting, apologies if nothing interesting was added (or taken away!)


\end{sffamily}\end{footnotesize}
